\item Un cascarón esférico tiene un radio externo de $2.60 \text{ cm}$ y uno interno de $a$. La pared del cascarón tiene grosor uniforme y está hecho de un material cuya densidad es de $4.70 \text{ g/cm}^3$. El espacio interior del cascarón está lleno con un líquido que tiene una densidad de $1.23 \text{ g/cm}^3$. (a) Encuentre la masa $m$ de la esfera, incluidos sus contenidos, como función de $a$. (b) ¿Para qué valor de $a$ tiene $m$ su máximo valor posible? (c) ¿Cuál es esta masa máxima? (d) Explique si el valor de la parte (c) concuerda con el resultado de un cálculo directo de la masa de una esfera sólida de densidad uniforme hecha del mismo material que el cascarón? (e) ¿Qué pasaría si? En el inciso (a), ¿la respuesta cambiaría si la pared interior del cascarón no fuese concéntrica con la pared exterior?
